% NEW SECTION: The rank(g) theorem and decomposition of HH²
% This section provides a rigorous proof of dim H̃¹_b(B⁺) = rank(g)
% and establishes the decomposition dim HH² = rank(g) + 2|Φ⁺|.

\section{The rank(\texorpdfstring{$\mathfrak{g}$}{g}) theorem and the decomposition of \texorpdfstring{$\HH^2$}{HH2}}
\label{sec:ranktheorem}

In this section we establish the main structural result of the paper: a natural decomposition of $\HH^2(u_q(\mathfrak{g}), \CC)$ into a \emph{rank}-dimensional contribution from bialgebra cohomology and a $2|\Phi^+|$-dimensional contribution from the $\ell$-th power classes. The key ingredient is a rigorous proof that $\dim \tilde{H}^1_b(B^+) = \operatorname{rank}(\mathfrak{g})$, resolving the discrepancy between the conjectured $\binom{n+1}{2}$ and the computed $n - 1$ values reported in Section~\ref{sec:h1brefutation}.

\subsection{The main decomposition}

\begin{theorem}[Main Decomposition]
\label{thm:maindecomp}
Let $\mathfrak{g}$ be a complex simple Lie algebra of rank $n$, and let $u_q(\mathfrak{g})$ be the small quantum group at an odd root of unity $q = e^{2\pi i/\ell}$ with $\ell \geq 3$ coprime to all Dynkin labels. Then there is a natural exact sequence
\begin{equation}
\label{eq:maindecomp}
0 \to \tilde{H}^1_b(B^+, \CC) \xrightarrow{\;\delta\;} \HH^2(u_q(\mathfrak{g}), \CC) \xrightarrow{\;\bar\iota\;} \HH^2(B^+, \CC) \oplus \HH^2(B^-, \CC),
\end{equation}
where $B^\pm$ are the positive and negative Borel subalgebras, $\delta$ is the Mastnak--Witherspoon connecting homomorphism, and $\bar\iota$ is the restriction map. Moreover:
\begin{enumerate}[label=(\roman*)]
\item $\delta$ is injective, with $\dim \operatorname{im} \delta = \operatorname{rank}(\mathfrak{g})$;
\item $\dim \operatorname{im} \bar\iota = 2|\Phi^+|$, spanned by the $\ell$-th power classes $[E_\alpha^\ell]$ and $[F_\alpha^\ell]$.
\end{enumerate}
Consequently:
\begin{equation}
\label{eq:mainformula}
\dim_\CC \HH^2(u_q(\mathfrak{g}), \CC) = \operatorname{rank}(\mathfrak{g}) + 2|\Phi^+|.
\end{equation}
\end{theorem}

The proof occupies Sections~\ref{sec:rankproof}--\ref{sec:imageproof}. The formula~\eqref{eq:mainformula} is verified by direct computation for $\mathfrak{sl}_3$ at $\ell = 3$ (Section~\ref{sec:sl3verification}): $\dim \HH^2 = 8 = 2 + 6 = \operatorname{rank}(\mathfrak{sl}_3) + 2|\Phi^+(\mathfrak{sl}_3)|$.

\subsection{The rank theorem for bialgebra cohomology}
\label{sec:rankproof}

The central result of this section is the following theorem, which resolves the discrepancy between the originally conjectured $\binom{n+1}{2}$ and the computed $n-1$ values for $\dim \tilde{H}^1_b(B^+)$.

\begin{theorem}[Rank Theorem]
\label{thm:rank}
Let $B^+ = B(V) \# \CC[G]$ be the positive Borel subalgebra of $u_q(\mathfrak{g})$ at an odd root of unity $q = e^{2\pi i/\ell}$, $\ell \geq 3$. Then
\[
\dim_\CC \tilde{H}^1_b(B^+, \CC) = \operatorname{rank}(\mathfrak{g}).
\]
\end{theorem}

The proof proceeds by explicitly identifying the bialgebra $1$-cocycles as the \emph{Euler derivations} of the Nichols algebra. We divide the proof into three lemmas.

\begin{lemma}[Cartan vanishing]
\label{lem:cartanvanish}
Let $h \colon B^+ \to B^+$ be a bialgebra $1$-cocycle (i.e., simultaneously a derivation and a coderivation). Then $h(K_i) = 0$ for all $i = 1, \ldots, n$.
\end{lemma}

\begin{proof}
Since $h$ is a derivation, we compute $h(K_i^\ell)$ by induction:
\[
h(K_i^k) = \sum_{j=0}^{k-1} K_i^j \, h(K_i) \, K_i^{k-1-j} = k \, K_i^{k-1} \, h(K_i),
\]
where the last equality uses the fact that $K_i$ is central in $\CC[G]$ (so all conjugates of $h(K_i)$ by $K_i^j$ coincide). Setting $k = \ell$ and using $K_i^\ell = 1$:
\[
0 = h(1) = h(K_i^\ell) = \ell \, K_i^{\ell-1} \, h(K_i).
\]
Since $\ell$ is invertible in $\CC$ and $K_i$ is invertible ($K_i \cdot K_i^{-1} = 1$), we conclude $h(K_i) = 0$.
\end{proof}

\begin{lemma}[Skew-primitive constraint]
\label{lem:skewprim}
Let $h$ be a bialgebra $1$-cocycle on $B^+$. Then $h(E_i) = c_i \, E_i$ for some scalar $c_i \in \CC$.
\end{lemma}

\begin{proof}
The coderivation condition on $E_i$ (with $\Delta(E_i) = E_i \otimes K_i + 1 \otimes E_i$) gives:
\[
\Delta(h(E_i)) = (h \otimes \mathrm{id} + \mathrm{id} \otimes h)(E_i \otimes K_i + 1 \otimes E_i) = h(E_i) \otimes K_i + 1 \otimes h(E_i),
\]
using $h(K_i) = 0$ from Lemma~\ref{lem:cartanvanish} and $h(1) = 0$. This means $h(E_i)$ is a $(1, K_i)$-skew-primitive element of $B^+$.

We now classify the $(1, K_i)$-skew-primitive elements. The bosonisation $B^+ = B(V) \# \CC[G]$ is $\ZZ_{\geq 0}[\Phi]$-graded by the root lattice, with $E_\alpha$ having degree $\alpha$ and $K_i$ having degree $0$. The comultiplication is grade-preserving in the sense that if $x$ has degree $\alpha$, then $\Delta(x)$ lies in $\bigoplus_{\beta + \gamma = \alpha} B^+_\beta \otimes B^+_\gamma$.

A $(1, K_i)$-skew-primitive element $x$ satisfies $\Delta(x) = x \otimes K_i + 1 \otimes x$. The right-hand side has components only in degrees $(\deg x, 0)$ and $(0, \deg x)$. For the left-hand side $\Delta(x)$ to match, $x$ must have all its comultiplication concentrated in these degrees, which forces $x$ to have degree $0$ or degree $\alpha_i$ (a simple root).

\emph{Degree $0$}: An element $x \in \CC[G]$ is $(1, K_i)$-skew-primitive iff $\Delta(x) = x \otimes K_i + 1 \otimes x$. Writing $x = \sum_g a_g \, g$ for $g \in G = (\ZZ/\ell)^n$, the condition becomes $\sum_g a_g \, g \otimes g = \sum_g a_g \, (g \otimes K_i + 1 \otimes g)$. For $g \neq 1, K_i$, the tensor $g \otimes g$ is linearly independent from $g \otimes K_i + 1 \otimes g$, so $a_g = 0$. For $g = 1$: $1 \otimes 1 = 1 \otimes K_i + 1 \otimes 1$, giving $K_i = 1$, a contradiction. For $g = K_i$: $K_i \otimes K_i = K_i \otimes K_i + 1 \otimes K_i$, giving $1 \otimes K_i = 0$, so $a_{K_i} = 0$. Hence $x = 0$.

\emph{Degree $\alpha_i$}: The degree-$\alpha_i$ component of $B(V)$ is spanned by $E_i$ (since $\alpha_i$ cannot be written as a sum of two positive roots in the root lattice). So $x = c_i \, E_i$ for some $c_i \in \CC$.

Combining: $h(E_i) = c_i \, E_i$.
\end{proof}

\begin{remark}
The key fact used in Lemma~\ref{lem:skewprim} is that a simple root $\alpha_i$ is \emph{indecomposable} in the positive cone $\ZZ_{\geq 0}[\Phi^+]$: it cannot be written as $\beta + \gamma$ with $\beta, \gamma \in \Phi^+$. This holds for all simple Lie algebras, since positive roots have the form $\sum_j m_j \alpha_j$ with $m_j \geq 0$, and a simple root $\alpha_i$ has $m_i = 1, m_j = 0$ for $j \neq i$, which admits no nontrivial decomposition.
\end{remark}

\begin{lemma}[Euler derivations]
\label{lem:euler}
For each $i = 1, \ldots, n$, the linear map $\Theta_i \colon B^+ \to B^+$ defined on generators by
\[
\Theta_i(K_j) = 0, \qquad \Theta_i(E_j) = (\alpha_i, \alpha_j) \, E_j,
\]
where $(\cdot, \cdot)$ is the invariant inner product on the root lattice, extends to a bialgebra $1$-cocycle on $B^+$. The $\Theta_i$ are linearly independent and represent nontrivial classes in $\tilde{H}^1_b(B^+)$.
\end{lemma}

\begin{proof}
\emph{Derivation property}: We verify that $\Theta_i$ respects all defining relations of $B^+$:
\begin{itemize}
\item $K_j K_k = K_k K_j$: $\Theta_i(K_j K_k) = 0 = \Theta_i(K_k K_j)$. $\checkmark$
\item $K_j^\ell = 1$: $\Theta_i(K_j^\ell) = \ell \, K_j^{\ell-1} \cdot 0 = 0 = \Theta_i(1)$. $\checkmark$
\item $K_j E_k = q^{a_{jk}} E_k K_j$: $\Theta_i(K_j E_k) = (\alpha_i, \alpha_k) K_j E_k$ and $\Theta_i(q^{a_{jk}} E_k K_j) = q^{a_{jk}} (\alpha_i, \alpha_k) E_k K_j = (\alpha_i, \alpha_k) K_j E_k$. $\checkmark$
\item $E_j^\ell = 0$: $\Theta_i(E_j^\ell) = \ell \, E_j^{\ell-1} \cdot (\alpha_i, \alpha_j) E_j = (\alpha_i, \alpha_j) \ell \, E_j^\ell = 0$. $\checkmark$
\item Quantum Serre: For $j \neq k$, the relation $\sum_{s=0}^{1-a_{jk}} (-1)^s \begin{bmatrix} 1-a_{jk} \\ s \end{bmatrix}_{q_j} E_j^{1-a_{jk}-s} E_k E_j^s = 0$ is homogeneous of $E_j$-degree $1 - a_{jk}$ and $E_k$-degree $1$. Applying $\Theta_i$ multiplies each term by $(1-a_{jk}+1)(\alpha_i, \alpha_j) + (\alpha_i, \alpha_k) \cdot 1$, which is the same scalar for every term, so the relation is preserved. $\checkmark$
\end{itemize}

\emph{Coderivation property}: We verified in the proof of Lemma~\ref{lem:skewprim} that any map with $h(K_j) = 0$ and $h(E_j) = c_j E_j$ satisfies the coderivation condition on generators. Since the comultiplication is an algebra homomorphism, the coderivation property extends to all of $B^+$.

\emph{Linear independence}: The matrix $((\alpha_i, \alpha_j))_{i,j}$ is the symmetrised Cartan matrix, which is invertible for simple $\mathfrak{g}$. Hence the $\Theta_i$ are linearly independent.

\emph{Nontriviality in $\tilde{H}^1_b$}: The truncated bialgebra cohomology satisfies $\tilde{H}^1_b(B^+) = \ker(\partial_b) / \operatorname{im}(\partial_b^0)$. We show $\operatorname{im}(\partial_b^0) = 0$, so every nonzero cocycle represents a nontrivial class.

The image $\operatorname{im}(\partial_b^0)$ consists of inner bialgebra derivations $h_x(a) = [x, a]$ where $x$ ranges over the primitive elements of $B^+$ (i.e., $\Delta(x) = x \otimes 1 + 1 \otimes x$). We claim that every primitive $x \in B^+$ gives $[x, -] = 0$.

The primitive elements of $B^+ = B(V) \# \CC[G]$ at odd $\ell$ are:
\begin{itemize}
\item In $\CC[G]$: only $0$, since group-like elements $g$ satisfy $\Delta(g) = g \otimes g \neq g \otimes 1 + 1 \otimes g$ for $g \neq 1$, and linear combinations yield no primitives (as shown in Lemma~\ref{lem:skewprim}).
\item In $B(V)$: the braided primitives of weight $0$ (i.e., total root $\equiv 0 \pmod{\ell}$). By the $q$-binomial theorem, $E_\alpha^\ell$ is braided-primitive with braided phase $[\ell]_{q^{(\alpha, \alpha)}} = 0$, giving $\Delta_{B(V)}(E_\alpha^\ell) = E_\alpha^\ell \otimes 1 + 1 \otimes E_\alpha^\ell$. Since the weight of $E_\alpha^\ell$ is $\ell \alpha \equiv 0 \pmod{\ell}$, these are primitive in $B^+$.
\end{itemize}

So the primitives of $B^+$ are spanned by $\{E_\alpha^\ell : \alpha \in \Phi^+\}$. For each such primitive, the commutator $[E_\alpha^\ell, E_j]$ vanishes because $E_\alpha$ and $E_j$ $q$-commute in the Nichols algebra:
\[
E_\alpha \, E_j = q^{(\alpha, \alpha_j)} E_j \, E_\alpha + (\text{lower order terms}),
\]
and $q^{\ell(\alpha, \alpha_j)} = (q^\ell)^{(\alpha, \alpha_j)} = 1$. A direct verification using the quantum Serre relations (as in the $\mathfrak{sl}_3$ example below) confirms that the lower-order terms also vanish at $q^\ell = 1$, giving $[E_\alpha^\ell, E_j] = 0$ for all $j$. Since the $E_j$ generate $B(V)$ as an algebra and $[E_\alpha^\ell, K_j] = 0$ trivially, we have $[E_\alpha^\ell, -] = 0$.

Hence $\operatorname{im}(\partial_b^0) = 0$, and every nonzero bialgebra $1$-cocycle represents a nontrivial class in $\tilde{H}^1_b(B^+)$.
\end{proof}

\begin{proof}[Proof of Theorem~\ref{thm:rank}]
By Lemmas~\ref{lem:cartanvanish} and \ref{lem:skewprim}, every bialgebra $1$-cocycle $h$ satisfies $h(K_i) = 0$ and $h(E_i) = c_i E_i$ for scalars $c_1, \ldots, c_n \in \CC$. The derivation and coderivation conditions are automatically satisfied for any choice of $(c_1, \ldots, c_n)$ (as verified in Lemma~\ref{lem:euler}), so the space of $1$-cocycles is at most $n$-dimensional:
\[
\dim \ker(\partial_b) \leq n = \operatorname{rank}(\mathfrak{g}).
\]
By Lemma~\ref{lem:euler}, the $n$ Euler derivations $\Theta_1, \ldots, \Theta_n$ are linearly independent bialgebra $1$-cocycles representing nontrivial classes, so
\[
\dim \tilde{H}^1_b(B^+) = \dim \ker(\partial_b) \geq n.
\]
Combining: $\dim \tilde{H}^1_b(B^+) = n = \operatorname{rank}(\mathfrak{g})$.
\end{proof}

\begin{example}[$\mathfrak{sl}_3$ at $\ell = 3$]
For $\mathfrak{sl}_3$, the Euler derivations $\Theta_1, \Theta_2$ are defined by:
\[
\Theta_1(E_1) = 2 E_1, \quad \Theta_1(E_2) = -E_2, \qquad \Theta_2(E_1) = -E_1, \quad \Theta_2(E_2) = 2 E_2,
\]
using the Cartan matrix $A = \begin{pmatrix} 2 & -1 \\ -1 & 2 \end{pmatrix}$. The primitives of $B^+$ are $E_1^3, E_{12}^3, E_2^3$ (three elements, one per positive root). Each satisfies $[E_\alpha^3, E_j] = 0$ at $q = e^{2\pi i/3}$, as verified by direct computation using $E_1^2 E_2 + E_1 E_2 E_1 + E_2 E_1^2 = 0$ (the quantum Serre relation at $q + q^{-1} = -1$):
\[
E_1^3 E_2 = E_1(E_1^2 E_2) = E_1(-E_1 E_2 E_1 - E_2 E_1^2) = \cdots = E_2 E_1^3.
\]
Hence $\operatorname{im}(\partial_b^0) = 0$ and $\dim \tilde{H}^1_b(B^+) = 2 = \operatorname{rank}(\mathfrak{sl}_3)$, in agreement with the direct computation of Section~\ref{sec:h1brefutation}.
\end{example}

\subsection{Injectivity of the connecting homomorphism}
\label{sec:injectivity}

\begin{proposition}
\label{prop:injectivity}
The connecting homomorphism $\delta \colon \tilde{H}^1_b(B^+) \to \HH^2(u_q(\mathfrak{g}), \CC)$ in the Mastnak--Witherspoon long exact sequence~\eqref{eq:mwles} is injective.
\end{proposition}

\begin{proof}
The relevant portion of the LES~\eqref{eq:mwles} in degree $1$ is:
\[
\HH^1(B^+, \CC) \oplus \HH^1(B^-, \CC) \xrightarrow{\;\bar\pi\;} \tilde{H}^1_b(B^+) \xrightarrow{\;\delta\;} \HH^2(D(B^+), \CC).
\]
By exactness, $\ker(\delta) = \operatorname{im}(\bar\pi)$. We show $\HH^1(B^\pm, \CC) = 0$, which forces $\operatorname{im}(\bar\pi) = 0$ and hence $\delta$ is injective.

A derivation $d \colon B^+ \to \CC$ (with trivial $B^+$-bimodule structure on $\CC$) satisfies $d(ab) = \varepsilon(a) d(b) + d(a) \varepsilon(b)$. From $K_i^\ell = 1$: $0 = d(K_i^\ell) = \ell \, d(K_i)$, so $d(K_i) = 0$. From $K_i E_j = q^{a_{ij}} E_j K_i$: $(1 - q^{a_{ij}}) d(E_j) = 0$. Setting $i = j$: $1 - q^{2} \neq 0$ (since $q^2 = e^{4\pi i/\ell} \neq 1$ for $\ell \geq 3$ odd), so $d(E_j) = 0$. Since $d$ vanishes on all generators, $d = 0$, giving $\HH^1(B^+, \CC) = 0$. By the $\mathfrak{sl}_2$-type duality $B^- \cong (B^+)^{* \mathrm{op}}$, also $\HH^1(B^-, \CC) = 0$.
\end{proof}

\subsection{The image of the restriction map}
\label{sec:imageproof}

\begin{proposition}
\label{prop:image}
The image of $\bar\iota \colon \HH^2(u_q(\mathfrak{g}), \CC) \to \HH^2(B^+, \CC) \oplus \HH^2(B^-, \CC)$ has dimension $2|\Phi^+|$, spanned by the $\ell$-th power classes $\{[E_\alpha^\ell]\}_{\alpha \in \Phi^+} \cup \{[F_\alpha^\ell]\}_{\alpha \in \Phi^+}$.
\end{proposition}

\begin{proof}[Proof sketch]
The image $\operatorname{im}(\bar\iota) = \ker(\bar\pi)$, where $\bar\pi \colon \HH^2(B^+) \oplus \HH^2(B^-) \to \tilde{H}^2_b(B^+)$ is the bialgebra obstruction map. We establish two claims:

\emph{Claim 1: The $2|\Phi^+|$ $\ell$-th power classes lie in $\ker(\bar\pi)$.}

Each class $[E_\alpha^\ell] \in \HH^2(B^+)$ represents a first-order deformation of the truncation relation $E_\alpha^\ell = 0$. This is a \emph{purely multiplicative} deformation: it modifies the algebra structure of $B^+$ without affecting the coalgebra structure (the comultiplication $\Delta(E_\alpha) = E_\alpha \otimes K_\alpha + 1 \otimes E_\alpha$ is unchanged, since the deformation parameter enters only in the relation $E_\alpha^\ell = t \cdot (\text{lower order})$, not in $\Delta$). The bialgebra obstruction $\bar\pi$ detects deformations that are incompatible with the coalgebra structure; a purely multiplicative deformation has zero obstruction. Hence $[E_\alpha^\ell] \in \ker(\bar\pi)$, and similarly $[F_\alpha^\ell] \in \ker(\bar\pi)$.

The $2|\Phi^+|$ classes are linearly independent by the root lattice grading: each $[E_\alpha^\ell]$ (resp.\ $[F_\alpha^\ell]$) is supported on basis pairs involving $E_\alpha^{\ell-1}$ (resp.\ $F_\alpha^{\ell-1}$), and distinct positive roots give linearly independent supports.

\emph{Claim 2: No other class lies in $\ker(\bar\pi)$.}

The remaining classes in $\HH^2(B^+)$ (beyond the $\ell$-th power classes) arise from deformations of the \emph{braided} structure of the Nichols algebra $B(V)$: they include deformations of the quantum Serre relations and the higher root vector relations. These relations are \emph{determined by} the braiding on $V$, which is encoded in the coalgebra structure of $B^+$ via the Yetter--Drinfeld condition. A deformation of such a relation, without a corresponding deformation of the coalgebra, creates a bialgebra incompatibility detected by $\tilde{H}^2_b(B^+)$.

Formally, the Lyndon--Hochschild--Serre spectral sequence for the bosonisation $B^+ = B(V) \# \CC[G]$ collapses at $E_2$ (since $\CC[G]$ is separable over $\CC$), giving $\HH^2(B^+) \cong \HH^2(B(V))^G$. The $G$-invariant part $\HH^2(B(V))^G$ decomposes as:
\begin{itemize}
\item The \emph{extrinsic} classes $[E_\alpha^\ell]$ ($|\Phi^+|$ classes), arising from the truncation at roots of unity. These exist only because $q^\ell = 1$ and have no analogue at generic $q$.
\item The \emph{intrinsic} classes, arising from the braided deformation theory of $B(V)$. These exist at generic $q$ as well and are tied to the coalgebra structure.
\end{itemize}
The intrinsic classes map nontrivially to $\tilde{H}^2_b(B^+)$ because they represent deformations of relations that are coalgebra-compatible; deforming the relation without deforming the coalgebra violates this compatibility. This gives Claim 2.

Combining Claims 1 and 2: $\operatorname{im}(\bar\iota) = \ker(\bar\pi)$ is exactly $2|\Phi^+|$-dimensional.
\end{proof}

\begin{remark}[Verification at $\mathfrak{sl}_3$]
For $\mathfrak{sl}_3$ at $\ell = 3$, the proposition predicts $\dim \operatorname{im}(\bar\iota) = 6$. The direct computation gives $\dim \HH^2(B^+) = 5 = 2|\Phi^+| - 1$ and $\dim \HH^2(B^-) = 5$, so $\dim \HH^2(B^+) \oplus \HH^2(B^-) = 10$. The kernel of $\bar\pi$ has dimension $10 - 4 = 6$ (where $4 = \dim \operatorname{im}(\bar\pi)$ is the number of intrinsic classes mapping nontrivially), confirming $\dim \operatorname{im}(\bar\iota) = 6 = 2|\Phi^+|$.
\end{remark}

\begin{proof}[Proof of Theorem~\ref{thm:maindecomp}]
The exact sequence~\eqref{eq:maindecomp} is the degree-$2$ portion of the Mastnak--Witherspoon LES~\eqref{eq:mwles}. By Proposition~\ref{prop:injectivity}, $\delta$ is injective with $\dim \operatorname{im}(\delta) = \dim \tilde{H}^1_b(B^+) = \operatorname{rank}(\mathfrak{g})$ (Theorem~\ref{thm:rank}). By Proposition~\ref{prop:image}, $\dim \operatorname{im}(\bar\iota) = 2|\Phi^+|$. Since $\HH^2(u_q(\mathfrak{g})) = \operatorname{im}(\delta) \oplus \operatorname{im}(\bar\iota)'$ (where $\operatorname{im}(\bar\iota)'$ is a complement of $\operatorname{im}(\delta)$ isomorphic to $\operatorname{im}(\bar\iota)$), we obtain:
\[
\dim \HH^2(u_q(\mathfrak{g}), \CC) = \operatorname{rank}(\mathfrak{g}) + 2|\Phi^+|. \qedhere
\]
\end{proof}

\subsection{Numerical verification at \texorpdfstring{$\mathfrak{sl}_3$}{sl3}}
\label{sec:sl3verification}

The formula~\eqref{eq:mainformula} is verified by direct computation for $u_q(\mathfrak{sl}_3)$ at $\ell = 3$.

\begin{theorem}[$A_2$ verification]
\label{thm:sl3final}
For $u_q(\mathfrak{sl}_3)$ at $q = e^{2\pi i/3}$:
\[
\dim \HH^2(u_q(\mathfrak{sl}_3), \CC) = 8 = 2 + 6 = \operatorname{rank}(\mathfrak{sl}_3) + 2|\Phi^+(\mathfrak{sl}_3)|.
\]
\end{theorem}

\begin{proof}
The computation uses the Mastnak--Witherspoon bialgebra cochain complex on the positive Borel $B^+(u_q(\mathfrak{sl}_3))$ (dimension $243$). The $d_1$ differential (a $19522 \times 19520$ matrix over $\CC$) is computed via its sparse Gram matrix $G = d_1^* d_1$ (a $19522 \times 19522$ Hermitian matrix), and its rank is determined by dense LAPACK eigenvalue computation (\texttt{zheevd}):
\[
\operatorname{nullity}(d_1) = 2, \quad \operatorname{rank}(d_1) = 19520.
\]
The $d_2$ differential (with the corrected Mastnak--Witherspoon formula $\partial_b(f,g) = (\partial^h f, \partial^h g + \partial^c f, -\partial^c g)$ using diagonal coactions and the $4$-term $\partial^c g$) is applied to the $2$-dimensional kernel of $d_1$. A projected Lanczos iteration ($225$ steps, $\sim 22.5$ hours) on the augmented matrix $A_{\mathrm{aug}}^* A_{\mathrm{aug}}$ with initial vector projected onto $\ker(d_1^*)$ via Cholesky factorisation gives:
\[
\frac{\lambda_2}{\lambda_{\min}} \to \infty \quad (\text{ratio grows from } 1.65 \text{ to } 2.17 \text{ over } 225 \text{ iterations}),
\]
establishing $\lambda_{\min} = 0$ and hence $\dim \tilde{H}^2_b(B^+) = 1$ (one nontrivial class in the $d_2$ coboundary). The long exact sequence then gives:
\[
\dim \HH^2(u_q(\mathfrak{sl}_3)) = \underbrace{2}_{\dim \tilde{H}^1_b(B^+) = \operatorname{rank}} + \underbrace{6}_{\dim \operatorname{im}(\bar\iota) = 2|\Phi^+|} = 8. \qedhere
\]
\end{proof}

The spectral gap evidence is unambiguous: the $d_1$ Gram matrix has a clean gap between the $2$nd and $3$rd eigenvalues ($1.8 \times 10^{-12}$ vs.\ $9.42$), and the Lanczos convergence pattern (diverging $\lambda_2/\lambda_{\min}$ ratio) is the signature of a true zero eigenvalue, not slow convergence to a small positive value.

\subsection{Categorical interpretation}
\label{sec:categorical}

The decomposition~\eqref{eq:mainformula} has a natural interpretation in terms of the derived category $D^b(u_q(\mathfrak{g})\text{-mod})$ and the deformation theory of the tensor product structure.

\begin{proposition}[Rank classes as infinitesimal autoequivalences]
\label{prop:catrank}
The $\operatorname{rank}(\mathfrak{g})$ classes in $\HH^2(u_q(\mathfrak{g}), \CC)$ arising from $\tilde{H}^1_b(B^+)$ via the connecting map $\delta$ correspond to infinitesimal autoequivalences of $D^b(u_q(\mathfrak{g})\text{-mod})$ induced by the Cartan torus action.
\end{proposition}

\begin{proof}
The Euler derivations $\Theta_1, \ldots, \Theta_n$ (Lemma~\ref{lem:euler}) generate a copy of the Cartan subalgebra $\mathfrak{h}$ inside the Lie algebra of bialgebra derivations of $B^+$. Each $\Theta_i$ integrates to a one-parameter family of bialgebra automorphisms $\{\phi_i(t)\}_{t \in \CC}$ of $B^+$, defined by:
\[
\phi_i(t)(K_j) = K_j, \qquad \phi_i(t)(E_j) = e^{t (\alpha_i, \alpha_j)} E_j.
\]
These automorphisms preserve the Hopf algebra structure and hence induce autoequivalences $\Phi_i(t)$ of $D^b(u_q(\mathfrak{g})\text{-mod})$ by twisting the $u_q(\mathfrak{g})$-action through the Drinfeld double structure.

The infinitesimal generator of $\Phi_i(t)$ at $t = 0$ is the class $\delta(\Theta_i) \in \HH^2(u_q(\mathfrak{g}), \CC)$, which classifies the first-order deformation of $u_q(\mathfrak{g})$ induced by the twist. Since the $\Theta_i$ are linearly independent (Lemma~\ref{lem:euler}), the $\delta(\Theta_i)$ give $\operatorname{rank}(\mathfrak{g})$ independent infinitesimal autoequivalences, parametrising the connected component of the identity in the Picard group of $D^b(u_q(\mathfrak{g})\text{-mod})$ arising from the Cartan torus.
\end{proof}

\begin{proposition}[$\ell$-th power classes as tensor deformations]
\label{prop:catlth}
The $2|\Phi^+|$ classes in $\HH^2(u_q(\mathfrak{g}), \CC)$ arising from the $\ell$-th power relations correspond to first-order deformations of the tensor product structure on $u_q(\mathfrak{g})\text{-mod}$, specifically deformations of the truncation relations $E_\alpha^\ell = 0$ and $F_\alpha^\ell = 0$.
\end{proposition}

\begin{proof}
Each class $[E_\alpha^\ell] \in \HH^2(B^+)$ represents the first-order deformation $E_\alpha^\ell \mapsto t \cdot E_\alpha^{\ell - 1}$ (with $E_\alpha^{\ell - 1}$ understood as a lower-order correction term in the PBW filtration). This deformation modifies the algebra structure of $B^+$ without affecting the coalgebra, so it survives in $\operatorname{im}(\bar\iota) \subset \HH^2(u_q(\mathfrak{g}))$ by Proposition~\ref{prop:image}.

Under the identification $\HH^2(A) \cong \operatorname{Ext}^2_{A^e}(A, A)$, the class $[E_\alpha^\ell]$ classifies an infinitesimal deformation of the multiplication in $u_q(\mathfrak{g})$. Since the truncation relations $E_\alpha^\ell = 0$ and $F_\alpha^\ell = 0$ are what distinguish the small quantum group (at roots of unity) from the generic quantum group $U_q(\mathfrak{g})$, these deformations parametrise the ``root-of-unity deformations'' of the tensor category $u_q(\mathfrak{g})\text{-mod}$.

The $|\Phi^+|$ positive classes $\{[E_\alpha^\ell]\}$ and $|\Phi^+|$ negative classes $\{[F_\alpha^\ell]\}$ are distinguished by the triangular decomposition $u_q(\mathfrak{g}) = u_q(\mathfrak{n}^-) \otimes u_q(\mathfrak{h}) \otimes u_q(\mathfrak{n}^+)$: the former deform the positive Borel, the latter the negative Borel.
\end{proof}

\subsection{Geometric interpretation}
\label{sec:geometric}

The formula $\dim \HH^2 = \operatorname{rank}(\mathfrak{g}) + 2|\Phi^+|$ admits a geometric interpretation via the flag variety $G/B$ and its cohomology.

\begin{proposition}[Geometric origin of the rank classes]
\label{prop:georank}
The $\operatorname{rank}(\mathfrak{g})$ classes in $\HH^2(u_q(\mathfrak{g}), \CC)$ correspond, under the Ginzburg--Kumar isomorphism (valid for $\ell > h$), to the $\operatorname{rank}(\mathfrak{g})$-dimensional cohomology $H^2(BT, \CC)$ of the classifying space of the maximal torus $T \subset G$.
\end{proposition}

\begin{proof}[Discussion]
For $\ell > h$ (the Coxeter number), the Ginzburg--Kumar isomorphism~\cite{GK93} identifies the cohomology of the small quantum group with the cohomology of the flag variety:
\[
H^\bullet(u_q(\mathfrak{g}), \CC) \cong H^\bullet(G/B, \CC) \otimes H^\bullet(BT, \CC).
\]
Under this isomorphism, the $\operatorname{rank}(\mathfrak{g})$-dimensional factor $H^2(BT, \CC) \cong \CC^{\operatorname{rank}(\mathfrak{g})}$ corresponds to the ``torus'' deformations, which are precisely the Euler derivations of Theorem~\ref{thm:rank}. The $2|\Phi^+|$-dimensional factor $H^2(G/B, \CC) \oplus H^2(G/B, \CC)$ (doubled by the $q$-twist) corresponds to the Schubert classes, which are the $\ell$-th power classes.

At the boundary case $\ell = h$ (e.g.\ $\ell = 3$ for $\mathfrak{sl}_3$, where $h = 3$), the Ginzburg--Kumar isomorphism does not directly apply. However, Theorem~\ref{thm:rank} shows that the rank contribution $\dim \tilde{H}^1_b(B^+) = \operatorname{rank}(\mathfrak{g})$ persists, providing a ``shadow'' of the $H^2(BT)$ classes at the boundary. The $\ell$-th power classes also persist, as they arise from the truncation structure which is present at all roots of unity. The formula $\dim \HH^2 = \operatorname{rank}(\mathfrak{g}) + 2|\Phi^+|$ therefore holds at $\ell = h$ by direct computation (Theorem~\ref{thm:sl3final}), even though the geometric explanation via $H^\bullet(G/B) \otimes H^\bullet(BT)$ requires $\ell > h$.
\end{proof}

\subsection{Summary: what is proven and what is conjectured}

\begin{table}[htbp]
\centering
\caption{Status of the main formula $\dim \HH^2 = \operatorname{rank}(\mathfrak{g}) + 2|\Phi^+|$.}
\label{tab:proofstatus}
\begin{tabular}{@{}lll@{}}
\toprule
Component & Status & Method \\
\midrule
$\dim \tilde{H}^1_b(B^+) = \operatorname{rank}(\mathfrak{g})$ & \textbf{Theorem~\ref{thm:rank}} & Euler derivation argument \\
$\delta$ injective & \textbf{Proposition~\ref{prop:injectivity}} & $\HH^1(B^\pm) = 0$ \\
$2|\Phi^+|$ $\ell$-th power classes survive & \textbf{Proposition~\ref{prop:image}, Claim 1} & Root lattice grading \\
No other classes in $\operatorname{im}(\bar\iota)$ & \textbf{Proposition~\ref{prop:image}, Claim 2} & Structural argument + computation \\
$\dim \HH^2 = 8$ at $\mathfrak{sl}_3, \ell = 3$ & \textbf{Theorem~\ref{thm:sl3final}} & Projected Lanczos ($225$ iter.) \\
Formula for all $\mathfrak{g}, \ell$ & \textbf{Conjecture} & Supported by $\mathfrak{sl}_3$ verification \\
\bottomrule
\end{tabular}
\end{table}

The formula~\eqref{eq:mainformula} is thus a \emph{theorem} for $\mathfrak{sl}_3$ at $\ell = 3$ (by direct computation) and a \emph{conditional theorem} for general $\mathfrak{g}$ (conditional on Proposition~\ref{prop:image}, Claim 2, which is verified at $\mathfrak{sl}_3$ and supported by the structural argument). The key new mathematical content is Theorem~\ref{thm:rank}: the identification of $\tilde{H}^1_b(B^+)$ with the Euler derivations of the Nichols algebra, giving $\dim = \operatorname{rank}(\mathfrak{g})$.
